\documentclass[10pt,a4paper]{amsart}
\usepackage[margin=19mm]{geometry}
\usepackage{amsmath,amssymb,mathtools}
\usepackage[scr=rsfs,cal=euler]{mathalfa}
\usepackage{xfrac}
\usepackage[unicode,hidelinks,bookmarks=false]{hyperref}
\newcommand{\ie}{i.e.\ }
\newcommand{\cf}{cf.\ }
\newcommand{\resp}{respectively\ }
\newcommand{\Subsection}{\subsection}
\providecommand{\nobreakdash}{\nobreak\hbox{-}}
\setlength{\emergencystretch}{2em}
\multlinegap0pt
\usepackage{amssymb}
\usepackage{xcolor}
\usepackage{mathtools}
\newtheorem{proposition}{Proposition}
\newcommand{\Par}[1]{\left(#1\right)}
\newcommand{\dd}{\textup{d}}
\newcommand{\inprod}[2]{\left\langle#1, #2\right\rangle}
\newcommand{\sfP}{\mathsf{P}}
\newcommand{\tpi}{\tilde{\pi}}
\newcommand{\va}{\mathbf{a}}
\newcommand{\vw}{\mathbf{w}}
\newcommand{\vy}{\mathbf{y}}
\newcommand{\vz}{\mathbf{z}}
\title{Functional Stochastic Localization}
\author{Anming Gu, Bobby Shi, Kevin Tian}
\date{Source: arXiv:2602.03999v2 (CC BY 4.0)}
\begin{document}
\maketitle

\noindent\emph{Source and licence.} Functional Stochastic Localization. Version arXiv:2602.03999v2 (CC BY 4.0), distributed by arXiv under the Creative Commons Attribution 4.0 International licence, http://creativecommons.org/licenses/by/4.0/, as recorded in the arXiv licence metadata for this exact version and shown on the abstract page https://arxiv.org/abs/2602.03999v2. Full licence notice: \texttt{licenses/arxiv-2602.03999-CC-BY-4.0.txt}.\par\smallskip

\noindent\textit{Editorial scope.} Counted unit: the article's proposition prop:renormalization-main (source0.tex lines 556--559). The article's complete original proof of it is delivered with it (source0.tex lines 560--592, one \texttt{proof} environment of 33 lines). No mathematics is written, repaired or re-derived: the statement and the proof are the article's own lines, in the article's own notation, and no retained line is substituted or re-flowed.  No further source-local context is needed: every cross-reference of every retained line resolves inside the two delivered spans.  Nothing was omitted from any retained span, so the package carries no omission record.  No retained line contains a citation, so the wrapper carries no bibliography and no bibliography entry.  No retained line contains a figure environment, an image-inclusion command or drawing code, so no image is delivered and the package declares no asset dependency.  Editorial formatting only, in addition: an AMS \texttt{amsart} preamble that restates the article's own declarations and macros used by the retained lines, namely the environments \texttt{proposition} on separate counters, together with the standard packages those lines need.  The retained environments print in this document as Proposition 1 in order of appearance, whereas the article prints the same items with the numbering of its own full document.  The complete native source of the article as distributed in the arXiv e-print for this exact version is bundled unchanged as \texttt{sources/193949/source0.tex}.
\begin{proposition}\label{prop:renormalization-main}
$\sfP_{\tau, \lambda}$ defines a discrete-time Markov semigroup.  Moreover, $\sfP_{0, \tau}f(\mathbf{0})=\mathbb{E}_{\tpi_{(\tau)}^Y}[f]$ and
\[\tpi_{(\tau)}^Y(\vy)=\exp\Par{-V_\tau(\vy)-\varphi^{*\tau}(\vy)}.\]
\end{proposition}

\begin{proof}
Using the fact that $\exp\Par{-\varphi^{*0}}=\delta_{\mathbf{0}}$ with respect to test functions, 
\[\sfP_{\tau, \tau}f(\va)=\exp(V_\tau(\va))\int f(\va+\vw)\exp\Par{-V_\tau(\va+\vw)-\varphi^{*0}(\vw)}\,\dd\vw=f(\va).\]
Let $\tau\leq \lambda\leq \sigma$.  First, define
\[
    g(\va):=\sfP_{\lambda, \sigma}f(\va)
    =\exp\Par{V_\lambda(\va)}\int f(\va+\vw)\exp\Par{-V_\sigma(\va+\vw)-\varphi^{*(\sigma-\lambda)}(\vw)}\,\dd\vw.
\]
Then,
\begin{align*}
    \sfP_{\tau, \lambda}g(\va)
    &=\exp\Par{V_\tau(\va)}\int g(\va+\vy)\exp\Par{-V_\lambda(\va+\vy)-\varphi^{*(\lambda-\tau)}(\vy)}\,\dd\vy\\
    &=\exp\Par{V_\tau(\va)}\\
     &\quad  \cdot \int \left[ \exp\Par{V_\lambda(\va+\vy)}\int f(\va+\vy+\vw)\exp\Par{-V_\sigma(\va+\vy+\vw)-\varphi^{*(\sigma-\lambda)}(\vw)}\,\dd\vw  \right]\\
     &\quad \cdot \exp\Par{-V_\lambda(\va+\vy)-\varphi^{*(\lambda-\tau)}(\vy)}\,\dd\vy\\
     &=\exp\Par{V_\tau(\va)}\\
     &\quad \cdot \int\int f(\va+\vy+\vw) \exp\Par{-V_\sigma(\va+\vy+\vw)-\varphi^{*(\sigma-\lambda)}(\vw)-\varphi^{*(\lambda-\tau)}(\vy)}\,\dd\vw\,\dd\vy\\
     &=\exp\Par{V_\tau(\va)}\\
     &\quad \cdot \int\int f(\va+\vz) \exp\Par{-V_\sigma(\va+\vz)-\varphi^{*(\sigma-\lambda)}(\vw)-\varphi^{*(\lambda-\tau)}(\vz-\vw)}\,\dd\vw\,\dd\vz\\
     &=\exp\Par{V_\tau(\va)}\int f(\va+\vz)\exp\Par{-V_\sigma(\va+\vz)-\varphi^{*(\sigma-\tau)}(\vz)}\,\dd\vz\\
     &=\sfP_{\tau, \sigma}f(\va),
\end{align*}
which proves the first statement. For the second statement,
\begin{equation}\label{eq:tpiY-test-fns}
    \begin{aligned}
    \sfP_{0, \tau}f(\mathbf{0})
    &=\exp(V_0(\mathbf{0}))\int f(\vw)\exp\Par{-V_\tau(\vw)-\varphi^{*\tau}(\vw)}\,\dd\vw\\
    &=\int\int  f(\vw)\exp\Par{\inprod{\vw}{\vz}-\varphi^{*\tau}(\vw)-\tau \psi(\vz)}\pi(\dd\vz)\,\dd\vw\\
    &=\int f(\vw)\tpi_{(\tau)}^Y(\dd\vw),
\end{aligned}
\end{equation}
as claimed.  The explicit formula for $\tpi_{(\tau)}^Y$ can be seen from the first line of \eqref{eq:tpiY-test-fns} as the definition of integrating functions $f$ with respect to $\tpi_{(\tau)}^Y$.
\end{proof}

\end{document}
